\documentclass[11pt,a4paper]{article}
\usepackage[utf8]{inputenc}
\usepackage[T1]{fontenc}
\usepackage[brazil]{babel}
\usepackage{geometry}
\geometry{top=2cm,bottom=2cm,left=2cm,right=2cm}
\usepackage{xcolor}
\definecolor{ifscgreen}{RGB}{16,140,80}
\definecolor{ifscdark}{RGB}{15,23,42}
\usepackage{hyperref}
\hypersetup{colorlinks=true,linkcolor=ifscgreen,urlcolor=ifscgreen}
\usepackage{booktabs}
\usepackage{tabularx}
\usepackage{longtable}
\usepackage{fancyhdr}
\pagestyle{fancy}
\fancyhf{}
\lhead{\textcolor{ifscgreen}{\textbf{IFSC Câmpus Garopaba}} -- Auditoria Normativa de Acervo (CST Gestão Ambiental)}
\rfoot{\thepage}

\begin{document}

\begin{center}
    {\large \textbf{INSTITUTO FEDERAL DE SANTA CATARINA}}\\[2pt]
    {\normalsize \textbf{CÂMPUS GAROPABA -- CST EM GESTÃO AMBIENTAL}}\\[6pt]
    {\LARGE \textbf{\textcolor{ifscgreen}{RELATÓRIO OFICIAL DE AUDITORIA BIBLIOGRÁFICA}}}\\[4pt]
    {\large \textbf{Adequação Normativa de Acervo \& Demanda de Compras de Exemplares Físicos}}\\[8pt]
    \rule{\linewidth}{1.5pt}
\end{center}

\vspace{0.3cm}
\noindent \textbf{Para:} David (Bibliotecário-Documentalista)\\
\textbf{De:} Comissão de Reformulação do PPC CST em Gestão Ambiental\\
\textbf{Data:} 28 de Agosto de 2026\\
\textbf{Assunto:} Diagnóstico de Cobertura do Catálogo Sophia \& Quantitativos de Exemplares Físicos

\vspace{0.4cm}
\section*{1. Critérios Normativos Institucionais}
A auditoria aplicou rigorosamente os parâmetros oficiais do IFSC:
\begin{itemize}
    \item \textbf{Bibliografia Básica:} Mínimo de 2 títulos por UC; cada título deve contar com ao menos \textbf{3 exemplares físicos} no acervo local do câmpus.
    \item \textbf{Bibliografia Complementar:} Mínimo de 3 títulos por UC; cada título deve contar com ao menos \textbf{1 exemplar físico} no acervo local do câmpus.
\end{itemize}

\section*{2. Síntese Executiva dos Indicadores}

\begin{table}[h!]
\centering
\begin{tabularx}{\linewidth}{|X|c|X|}
\hline
\textbf{Indicador Normativo} & \textbf{Total} & \textbf{Observação / Diagnóstico} \\ \hline
Unidades Curriculares Auditadas & 33 UCs & 100\% da matriz curricular do CST Gestão Ambiental \\ \hline
Total de Referências Bibliográficas & 189 & 85 Básicas e 128 Complementares \\ \hline
Títulos Disponíveis no Acervo Sophia & 123 (65.1\%) & Obras com presença física confirmada no câmpus \\ \hline
Títulos Ausentes no Acervo Sophia & 66 (34.9\%) & Demanda prioritária para aquisição \\ \hline
Variações de Edição / Ano & 20 & Obras presentes com edição divergente da citada \\ \hline
\textbf{DEMANDA DE COMPRAS: Básica ($<$ 3 ex.)} & \textbf{+106 ex.} & \textbf{Exemplares para meta de 3 cópias na básica} \\ \hline
\textbf{DEMANDA DE COMPRAS: Complementar ($<$ 1 ex.)} & \textbf{+44 ex.} & \textbf{Exemplares para meta de 1 cópia na complementar} \\ \hline
\textbf{DEMANDA TOTAL DE COMPRAS} & \textbf{\textcolor{ifscgreen}{+150 EX.}} & \textbf{Total consolidado para conformidade 100\%} \\ \hline
\end{tabularx}
\end{table}

\newpage
\section*{3. Lista Prioritária de Aquisições: Bibliografia BÁSICA (+106 exemplares)}

\begin{longtable}{|p{1.2cm}|p{3.5cm}|p{6cm}|p{3.5cm}|c|}
\hline
\textbf{Sem.} & \textbf{Unidade Curricular} & \textbf{Título Referenciado no PPC} & \textbf{Autor} & \textbf{Déficit} \\ \hline
\endfirsthead
\hline
\textbf{Sem.} & \textbf{Unidade Curricular} & \textbf{Título Referenciado no PPC} & \textbf{Autor} & \textbf{Déficit} \\ \hline
\endhead
1º & ÉTICA AMBIENTAL & O princípio responsabilidade: ensaio de uma ética para a civilização tecnológica & JONAS, Hans & \textbf{+2 ex.} \\ \hline
1º & ÉTICA AMBIENTAL & O princípio vida: fundamentos para uma biologia filosófica & JONAS, Hans & \textbf{+3 ex.} \\ \hline
1º & ÉTICA AMBIENTAL & Ética ambiental & JUNGES, José Roque & \textbf{+3 ex.} \\ \hline
1º & ÉTICA AMBIENTAL & Nietzsche e a planta homem: sobre o perspectivismo vegetal & OLIVEIRA, Jelson & \textbf{+3 ex.} \\ \hline
1º & HISTÓRIA AMBIENTAL & A ferro e fogo: a história e a devastação da Mata Atlântica brasileira & DEAN, Warren & \textbf{+1 ex.} \\ \hline
1º & HISTÓRIA AMBIENTAL & Transformações da terra: para uma perspectiva agroecológica da história & WORSTER, Donald & \textbf{+3 ex.} \\ \hline
1º & LINGUAGEM E COMUNICAÇÃO I & Técnicas de comunicação escrita & BLIKSTEIN, Izidoro & \textbf{+2 ex.} \\ \hline
1º & METODOLOGIA DE PESQUISA & Metodologia do trabalho científico: procedimentos básicos, pesquisa bibliográfica, projeto e relatório, publicações e trabalhos científicos & MARCONI, Marina de Andrade; LAKATOS, Eva Maria & \textbf{+2 ex.} \\ \hline
1º & SEGURANÇA E SAÚDE OCUPACIONAL & Segurança do trabalho e gestão ambiental & BARBOSA FILHO, Antonio Nunes & \textbf{+1 ex.} \\ \hline
1º & SEGURANÇA E SAÚDE OCUPACIONAL & Segurança no trabalho e prevenção de acidentes: uma abordagem holística: segurança integrada à missão organizacional com produtividade, qualidade, preservação ambiental e desenvolvimento de pessoas & CARDELLA, Benedito & \textbf{+2 ex.} \\ \hline
2º & EDUCAÇÃO AMBIENTAL & Questão ambiental: questões de vida & LOUREIRO, Carlos Frederico & \textbf{+3 ex.} \\ \hline
2º & EDUCAÇÃO AMBIENTAL & O que é educação ambiental & REIGOTA, Marcos & \textbf{+1 ex.} \\ \hline
2º & FUNDAMENTOS DA ADMINISTRAÇÃO & Introdução à administração: da competitividade à sustentabilidade & DIAS, Reinaldo; ZAVAGLIA, Tércia; CASSAR, Maurício & \textbf{+2 ex.} \\ \hline
2º & FUNDAMENTOS DA ADMINISTRAÇÃO & Gestão ambiental: responsabilidade social e sustentabilidade & DIAS, Reinaldo & \textbf{+1 ex.} \\ \hline
2º & FUNDAMENTOS DA ADMINISTRAÇÃO & Administração: princípios e tendências & LACOMBE, F. J. M. ; HEILBORN, G. L. J & \textbf{+2 ex.} \\ \hline
2º & GEOMÁTICA & A interpretação de imagens aéreas: noções básicas e algumas aplicações nos campos profissionais & LOCH, Carlos & \textbf{+2 ex.} \\ \hline
2º & LINGUAGEM E COMUNICAÇÃO II & Aprenda a se comunicar com habilidade e clareza: 24 técnicas para tornar sua comunicação mais eficiente e seu dia-a-dia mais produtivo & ARREDONDO, Lani & \textbf{+2 ex.} \\ \hline
2º & MICROBIOLOGIA E TOXICOLOGIA AMBIENTAL & Microbiologia ambiental & ROCHA, Maria Carolina Vieira da & \textbf{+1 ex.} \\ \hline
2º & MICROBIOLOGIA E TOXICOLOGIA AMBIENTAL & Princípios de toxicologia ambiental: conceito e aplicações & SISINNO, Cristina Lúcia Silveira; OLIVEIRA-FILHO, Eduardo Cyrino & \textbf{+1 ex.} \\ \hline
3º & ECOLOGIA APLICADA & Fundamentos de ecologia & ODUM, E. P. ; BARRETT, G. W & \textbf{+3 ex.} \\ \hline
3º & GESTÃO DE ÁREAS PROTEGIDAS & Unidades de conservação: abordagens e características geográficas & GUERRA, Antonio José Teixeira; COELHO, Maria Célia Nunes (org.) & \textbf{+2 ex.} \\ \hline
3º & GESTÃO DE RECURSOS HÍDRICOS & Sistema Nacional de Informações sobre Recursos Hídricos: Conjuntura dos recursos hídricos no Brasil: Disponível em: https://www.snirh.gov.br/portal/snirh/centrais-de-conteudos/conjuntura-dos-recursos-hidricos Acesso em: 15 jul & AGÊNCIA NACIONAL DE ÁGUAS & \textbf{+3 ex.} \\ \hline
3º & GESTÃO DE RESÍDUOS SÓLIDOS & Política nacional, gestão e gerenciamento de resíduos sólidos & JARDIM, Arnaldo; YOSHIDA, Consuelo Yatsuda Moromizato; MACHADO FILHO, José Valverde & \textbf{+2 ex.} \\ \hline
3º & GESTÃO DE RESÍDUOS SÓLIDOS & Compostagem: nada se cria, nada se perde; tudo se transforma & MASSUKADO, Luciana Miyoko & \textbf{+1 ex.} \\ \hline
3º & PATRIMÔNIO CULTURAL NA GESTÃO AMBIENTAL & Patrimônio Cultural arqueológico: Diálogos, Reflexões e práticas & BASTOS, Rossano Lopes; SOUZA, Marise Campos de (Org.) & \textbf{+3 ex.} \\ \hline
3º & PATRIMÔNIO CULTURAL NA GESTÃO AMBIENTAL & Normas e gerenciamento do patrimônio arqueológico & BASTOS, R. L. ; SOUZA, Marise Campos de & \textbf{+3 ex.} \\ \hline
3º & TOPOGRAFIA APLICADA AO GEORREFERENCIAMENTO & Fundamentos de topografia & SARAIVA, Sérgio; TULER, Marcelo & \textbf{+3 ex.} \\ \hline
4º & AGROECOLOGIA APLICADA & Agroecologia: bases científicas para uma agricultura sustentável & ALTIERI, Miguel A & \textbf{+3 ex.} \\ \hline
4º & AGROECOLOGIA APLICADA & Agroecologia: alguns conceitos e princípios & CAPORAL, F. R. ; COSTABEBER, J. A & \textbf{+3 ex.} \\ \hline
4º & AGROECOLOGIA APLICADA & 19 Lições de Pedologia & LEPSCH, I. F & \textbf{+1 ex.} \\ \hline
4º & AGROECOLOGIA APLICADA & Conservação do solo & BERTONI, J. ; LOMBARDI NETO, F & \textbf{+3 ex.} \\ \hline
4º & FERRAMENTAS DE GESTÃO AMBIENTAL & Ferramentas de gestão ambiental: competitividade e sustentabilidade & PIMENTA, Handson Cláudio Dias; GOUVINHAS, Reidson Pereira & \textbf{+3 ex.} \\ \hline
4º & LEGISLAÇÃO AMBIENTAL & Direito ambiental brasileiro & MACHADO, Paulo Affonso Leme & \textbf{+1 ex.} \\ \hline
4º & LEGISLAÇÃO AMBIENTAL & Legislação Ambiental Comentada & TRENNEPOHL, Curt; TRENNEPOHL, Natascha; Terence & \textbf{+3 ex.} \\ \hline
4º & PROJETO DE TRABALHO DE CONCLUSÃO DE CURSO & Fundamentos de metodologia científica & LAKATOS, Eva Maria; MARCONI, Marina de Andrade & \textbf{+3 ex.} \\ \hline
4º & PROJETO DE TRABALHO DE CONCLUSÃO DE CURSO & Metodologia do trabalho científico & SEVERINO, Antônio Joaquim & \textbf{+3 ex.} \\ \hline
4º & TRATAMENTO DE ÁGUAS E EFLUENTES & Fundamentos de qualidade e tratamento de água & LIBÂNIO, Marcelo & \textbf{+1 ex.} \\ \hline
5º & AVALIAÇÃO DE IMPACTOS AMBIENTAIS & G; ALMEIDA, A & ALMEIDA, F. S. ; GARRIDO, F de S. R & \textbf{+3 ex.} \\ \hline
5º & EMPREENDEDORISMO & Inovação e empreendedorismo & BESSANT, John; TIDD, Joe; COSTA, Francisco Araújo da & \textbf{+1 ex.} \\ \hline
5º & EMPREENDEDORISMO & Business model generation = inovação em modelos de negócios : um manual para visionários, inovadores e revolucionários & OSTERWALDER, Alexander; PIGNEUR, Yves & \textbf{+2 ex.} \\ \hline
5º & LICENCIAMENTO AMBIENTAL & 190, de 8 de agosto de 2025 & BRASIL. Lei nº 15 & \textbf{+3 ex.} \\ \hline
1º & PLANEJAMENTO E GESTÃO DO TURISMO RESPONSÁVEL & São Paulo: Roca, 2005 & DESENVOLVIMENTO sustentável do turismo: uma compilação de boas práticas & \textbf{+2 ex.} \\ \hline
1º & PLANEJAMENTO E GESTÃO DO TURISMO RESPONSÁVEL & Gestão pública do turismo Brasil: teorias, metodologias e aplicações & PIMENTEL, Thiago Duarte; EMMENDOERFER, Magnus Luiz; TOMAZZONI, Edegar Luis (org.) & \textbf{+2 ex.} \\ \hline
5º & ESPANHOL APLICADO & Al día: curso de español para los negocios & PROST, Gisèle; NORIEGA FERNÁNDEZ, Alfredo & \textbf{+2 ex.} \\ \hline
5º & INGLÊS APLICADO & Gramática prática da língua inglesa: o inglês descomplicado & TORRES, Nelson & \textbf{+3 ex.} \\ \hline
5º & LIBRAS & Língua de sinais brasileira: estudos linguísticos & QUADROS, Ronice Müller de; KARNOPP, Lodenir Becker & \textbf{+1 ex.} \\ \hline
5º & LIBRAS & Introdução ao estudo da Libras & SILVA, Ronice Müller de Quadros, Rodrigo Nogueira Machado, Jair Barbosa da & \textbf{+3 ex.} \\ \hline
5º & LIBRAS & Aprenda a ver & WILCOX, Sherman; WILCOX, Phyllis Perrin & \textbf{+3 ex.} \\ \hline
\end{longtable}

\newpage
\section*{4. Lista de Aquisições: Bibliografia COMPLEMENTAR (+44 exemplares)}

\begin{longtable}{|p{1.2cm}|p{3.5cm}|p{6cm}|p{3.5cm}|c|}
\hline
\textbf{Sem.} & \textbf{Unidade Curricular} & \textbf{Título Referenciado no PPC} & \textbf{Autor} & \textbf{Déficit} \\ \hline
\endfirsthead
\hline
\textbf{Sem.} & \textbf{Unidade Curricular} & \textbf{Título Referenciado no PPC} & \textbf{Autor} & \textbf{Déficit} \\ \hline
\endhead
1º & ÉTICA AMBIENTAL & Fenomenologia da vida & FALABRETTI, Ericson; OLIVEIRA, Jelson (Orgs.) & +1 ex. \\ \hline
1º & ÉTICA AMBIENTAL & A nova ordem ecológica: a árvore, o animal e o homem & FERRY, Luc & +1 ex. \\ \hline
1º & ÉTICA AMBIENTAL & Gaia: alerta final & LOVELOCK, James & +1 ex. \\ \hline
1º & ÉTICA AMBIENTAL & Hans Jonas e a filosofia da vida & OLIVEIRA, J. R. ; MALTA, P. J & +1 ex. \\ \hline
1º & ÉTICA AMBIENTAL & Nietzsche, um proto-ambientalista? Poiesis: Revista de Filosofia, v.12, n.2, p.4-29, 2015 & OLIVEIRA, Jelson & +1 ex. \\ \hline
1º & HISTÓRIA AMBIENTAL & Primavera silenciosa & CARSON, Rachel & +1 ex. \\ \hline
1º & HISTÓRIA AMBIENTAL & Pensar como uma montanha & LEOPOLD, Aldo & +1 ex. \\ \hline
1º & HISTÓRIA AMBIENTAL & Fim do futuro? Manifesto ecológico brasileiro & LUTZENBERGER, José A & +1 ex. \\ \hline
1º & HISTÓRIA AMBIENTAL & Ecologia e política no Brasil & PÁDUA, José Augusto (org.) & +1 ex. \\ \hline
1º & HISTÓRIA AMBIENTAL & História ambiental: possibilidades e desafios & SEDREZ, Lise Fernanda & +1 ex. \\ \hline
1º & HISTÓRIA AMBIENTAL & Lucro: uma história ambiental & STOLL, Mark & +1 ex. \\ \hline
1º & METODOLOGIA DE PESQUISA & Trabalhos de pesquisa: diários de leitura para a revisão bibliográfica & LOUSADA, Eliane; ABREU-TARDELLI, Lília Santos (Org.) & +1 ex. \\ \hline
1º & SEGURANÇA E SAÚDE OCUPACIONAL & Normas Regulamentadoras & BRASIL. Ministério do trabalho & +1 ex. \\ \hline
2º & EDUCAÇÃO AMBIENTAL & Meio ambiente e consumismo & COIMBRA, José de Ávila Aguiar, 1930- (coord.) & +1 ex. \\ \hline
2º & GEOMÁTICA & In: XXV CONGRESSO BRASILEIRO DE CARTOGRAFIA, 2011, Curitiba & CASSOL, R. et al. Comparação entre câmeras fotográficas digitais com GPS para aplicação em geotagging & +1 ex. \\ \hline
3º & ECOLOGIA APLICADA & Manual de ecossistemas marinhos e costeiros para educadores & GERLING, C. et al (org.) & +1 ex. \\ \hline
3º & ECOLOGIA APLICADA & Curso de recuperação de áreas degradadas & BRASIL. Ministério da Agricultura, Pecuária e Abastecimento & +1 ex. \\ \hline
3º & ESTATÍSTICA APLICADA & Estatística aplicada às ciências humanas e ao turismo & RABAHY, W. A & +1 ex. \\ \hline
3º & GESTÃO DE RECURSOS HÍDRICOS & Brasília, 2003 & BANCO MUNDIAL. Estratégias de gerenciamento de recursos hídricos no Brasil: áreas de cooperação com o Banco Mundial & +1 ex. \\ \hline
3º & GESTÃO DE RECURSOS HÍDRICOS & Plano Nacional de Recursos Hídricos & BRASIL. Ministério do Meio Ambiente, Secretaria de Recursos Hídricos & +1 ex. \\ \hline
3º & PATRIMÔNIO CULTURAL NA GESTÃO AMBIENTAL & Patrimônio cultural plural & CAMPOS, Jussef Daibert Salomão de (org.) & +1 ex. \\ \hline
3º & PATRIMÔNIO CULTURAL NA GESTÃO AMBIENTAL & Paisagem cultural e patrimônio & RIBEIRO, Rafael Winter & +1 ex. \\ \hline
3º & TOPOGRAFIA APLICADA AO GEORREFERENCIAMENTO & Instituto Nacional de Colonização e Reforma Agrária & BRASIL. Ministério do Desenvolvimento Agrário & +1 ex. \\ \hline
3º & TOPOGRAFIA APLICADA AO GEORREFERENCIAMENTO & \& SEIXAS, A & SILVA, A. G. O. ; AZEVEDO, V. W. B & +1 ex. \\ \hline
3º & TOPOGRAFIA APLICADA AO GEORREFERENCIAMENTO & Fundamentos de topografia & VEIGA, L. A. K. ; ZANETTI, M. A. Z. ; FAGGION, P. L & +1 ex. \\ \hline
4º & AGROECOLOGIA APLICADA & Agroecologia: processos ecológicos em agricultura sustentável & GLIESSMAN, S. R & +1 ex. \\ \hline
4º & AGROECOLOGIA APLICADA & 5. ed & EMBRAPA. Sistema Brasileiro de Classificação de Solos & +1 ex. \\ \hline
4º & AGROECOLOGIA APLICADA & Plano Nacional de Agroecologia e Produção Orgânica (PLANAPO) & BRASIL. Ministério do Desenvolvimento Agrário e Agricultura Familiar & +1 ex. \\ \hline
4º & AGROECOLOGIA APLICADA & Plano Setorial de Mitigação e Adaptação às Mudanças Climáticas para a Consolidação de uma Economia de Baixa Emissão de Carbono na Agricultura (Plano ABC+) & BRASIL. Ministério da Agricultura, Pecuária e Abastecimento & +1 ex. \\ \hline
4º & AGROECOLOGIA APLICADA & Caminhos para o desenvolvimento sustentável & SACHS, I & +1 ex. \\ \hline
4º & AGROECOLOGIA APLICADA & Coleção Transição Agroecológica, Vol & EMBRAPA. Solos e agroecologia & +1 ex. \\ \hline
4º & ATIVIDADES DE EXTENSÃO & Pesquisa alienada e ensino alienante: o equívoco da extensão universitária & BOTOMÉ, Silvio Paulo & +1 ex. \\ \hline
4º & FERRAMENTAS DE GESTÃO AMBIENTAL & Indicadores de desempenho dos Sistemas de Gestão Ambiental (SGA): uma pesquisa teórica & CAMPOS, L. M. S. ; MELO, D. A & +1 ex. \\ \hline
4º & FERRAMENTAS DE GESTÃO AMBIENTAL & de F & MARTINS, M & +1 ex. \\ \hline
4º & PROJETO DE TRABALHO DE CONCLUSÃO DE CURSO & Projeto de pesquisa: métodos qualitativo, quantitativo e misto & CRESWELL, John W & +1 ex. \\ \hline
4º & PROJETO DE TRABALHO DE CONCLUSÃO DE CURSO & Métodos e técnicas de pesquisa social & GIL, Antonio Carlos & +1 ex. \\ \hline
4º & PROJETO DE TRABALHO DE CONCLUSÃO DE CURSO & Técnicas de pesquisa & MARCONI, Marina de Andrade; LAKATOS, Eva Maria & +1 ex. \\ \hline
4º & PROJETO DE TRABALHO DE CONCLUSÃO DE CURSO & Estudo de caso: planejamento e métodos & YIN, Robert K & +1 ex. \\ \hline
4º & PROJETO DE TRABALHO DE CONCLUSÃO DE CURSO & Metodologia do trabalho científico: métodos e técnicas da pesquisa e do trabalho acadêmico & PRODANOV, Cleber Cristiano; FREITAS, Ernani Cesar de & +1 ex. \\ \hline
5º & AVALIAÇÃO DE IMPACTOS AMBIENTAIS & Revista Monografias Ambientais, Santa Maria, RS, v.13, n.5, p.3821--3830, dez & CREMONEZ, F. E. et al. Avaliação de impacto ambiental: metodologias aplicadas no Brasil & +1 ex. \\ \hline
5º & AVALIAÇÃO DE IMPACTOS AMBIENTAIS & de; MEDEIROS, W & OLIVEIRA, F. F. G & +1 ex. \\ \hline
5º & LICENCIAMENTO AMBIENTAL & Cartilha de licenciamento ambiental & BRASIL. Tribunal de Contas da União & +1 ex. \\ \hline
5º & LIBRAS & 2. ed & MORAIS, Carlos E. L. ; PLINSKI, Rejane R. K. ; MARTINS, Gabriel P. T. C. ; et al. Libras & +1 ex. \\ \hline
5º & LIBRAS & Despertar do silêncio & VILHALVA, Shirley & +1 ex. \\ \hline
\end{longtable}

\end{document}
